\documentclass[12pt]{article}
\usepackage[utf8]{inputenc}
\usepackage[T1]{fontenc}
\usepackage[margin=1in]{geometry}
\usepackage{setspace}
\usepackage{xcolor}
\usepackage[colorlinks=true, urlcolor=blue!60!black, linkcolor=blue!60!black]{hyperref}
\setlength{\parindent}{2em}
\setlength{\parskip}{1.25em}
\renewcommand{\baselinestretch}{1.0}
\begin{document}
\begin{center}
{\LARGE \textit{Saenz v. Roe}, 526 U.S. 489 (1999)\par}
\vspace{0.8em}
{\large Case Brief\par}
\end{center}
\vspace{1em}

\section{Memory Jogger}

Reviving the Fourteenth Amendment's Privileges or Immunities Clause, the Supreme Court held 7--2 that California violated newly arrived residents' constitutional right to travel by capping their first-year welfare benefits at the level of their prior states of residence, and that Congress could not authorize states to do so.

\section{Detailed Case Facts}

California maintained one of the Nation's most generous public assistance programs, ranking sixth highest among the states in nominal welfare benefit levels (with forty-four states and the District of Columbia paying less), though eighteenth when housing costs were factored in (\href{https://www.law.cornell.edu/supremecourt/text/98-97}{Cornell LII}). Under Aid to Families with Dependent Children (AFDC)---later replaced by Temporary Assistance to Needy Families (TANF)---California provided cash assistance to an average of 2,645,814 persons per month at an annual cost of approximately \$2.9 billion (\href{https://www.law.cornell.edu/supremecourt/text/98-97}{Cornell LII}). For example, a family of two (a mother and one child) received \$456 per month in California compared to \$275 in neighboring Arizona, and a family of four received \$673 in California compared to \$144 in Mississippi (\href{https://www.law.cornell.edu/supremecourt/text/98-97}{Cornell LII}).

In 1992, seeking to reduce its welfare budget by an estimated \$10.9 million annually, the California Legislature enacted Cal. Welf. \& Inst. Code Ann. \S~11450.03 (\href{https://www.law.cornell.edu/supremecourt/text/98-97}{Cornell LII}). That statute limited the maximum welfare grant payable to any family residing in California for less than twelve consecutive months to the maximum amount the family would have received in its state of prior residence (\href{https://www.law.cornell.edu/supremecourt/text/98-97}{Cornell LII}). The cap applied regardless of whether the family had received welfare in its former state, regardless of its motive for moving to California, and even to lifelong Californians who briefly relocated elsewhere before returning; conversely, it did not apply to newcomers arriving from another country or from the five states with higher benefits (\href{https://www.law.cornell.edu/supremecourt/text/98-97}{Cornell LII}).

In 1996, Congress enacted the Personal Responsibility and Work Opportunity Reconciliation Act of 1996 (PRWORA), 110 Stat. 2105, replacing AFDC with TANF block grants and expressly authorizing any receiving state to ``apply to a family the rules (including benefit amounts) of the [TANF] program \ldots of another State if the family has moved to the State from the other State and has resided in the State for less than 12 months,'' 42 U.S.C. \S~604(c) (\href{https://www.law.cornell.edu/supremecourt/text/98-97}{Cornell LII}). Relying on \S~604(c), the California Department of Social Services issued an ``All County Letter'' announcing that enforcement of \S~11450.03 would begin on April 1, 1997 (\href{https://www.law.cornell.edu/supremecourt/text/98-97}{Cornell LII}). On that date, respondents Brenda Roe (who moved from Oklahoma) and Anna Doe (who moved from the District of Columbia)---both eligible for TANF benefits in jurisdictions with substantially lower grants than California---sued on behalf of a certified class of present and future applicants who would be denied full California benefits under the one-year durational residency cap (\href{https://www.law.cornell.edu/supremecourt/text/98-97}{Cornell LII}). After oral argument in the Supreme Court, Rita L. Saenz replaced Eloise Anderson as Director of the California Department of Social Services (\href{https://www.law.cornell.edu/supremecourt/text/98-97}{Cornell LII}).

\section{Procedural History}

\begin{itemize}
\item \textbf{Initial Challenge (\textit{Green v. Anderson}, 1992--1995):} After the Secretary of Health and Human Services issued a federal waiver approving \S~11450.03 in October 1992, three newly arrived California mothers (escaping abusive family circumstances in Louisiana, Oklahoma, and Colorado) sued in the U.S. District Court for the Eastern District of California. Judge David F. Levi preliminarily enjoined enforcement under \href{https://supreme.justia.com/cases/federal/us/394/618/}{\textit{Shapiro v. Thompson}}, 394 U.S. 618 (1969), and \href{https://supreme.justia.com/cases/federal/us/457/55/}{\textit{Zobel v. Williams}}, 457 U.S. 55 (1982), in \textit{Green v. Anderson}, 811 F. Supp. 516 (E.D. Cal. 1993), and the Ninth Circuit summarily affirmed, 26 F.3d 95 (9th Cir. 1994) (\href{https://www.law.cornell.edu/supremecourt/text/98-97}{Cornell LII}). The Supreme Court granted certiorari, 513 U.S. 922 (1994), but vacated and ordered dismissal as unripe in \textit{Anderson v. Green}, 513 U.S. 557 (1995) (per curiam), after the Ninth Circuit invalidated the Secretary's waiver in a separate case, \textit{Beno v. Shalala}, 30 F.3d 1057 (9th Cir. 1994) (\href{https://www.law.cornell.edu/supremecourt/text/98-97}{Cornell LII}).
\item \textbf{Post-PRWORA District Court Proceedings:} On April 1, 1997, following Congress's enactment of 42 U.S.C. \S~604(c), respondents filed this class action in the U.S. District Court for the Eastern District of California challenging both \S~11450.03 and PRWORA's authorization. Judge Levi certified the class and again preliminarily enjoined enforcement, holding that PRWORA's permissive authorization did not alter the constitutional analysis (\textit{Roe v. Anderson}, 966 F. Supp. 977 (E.D. Cal. 1997)) (\href{https://www.law.cornell.edu/supremecourt/text/98-97}{Cornell LII}; \href{https://journals.law.harvard.edu/jlpp/wp-content/uploads/sites/90/2026/03/30_23HarvJLPubPoly2951999-2000.pdf}{Harvard JLPP}).
\item \textbf{Court of Appeals:} The U.S. Court of Appeals for the Ninth Circuit affirmed the preliminary injunction in \textit{Roe v. Anderson}, 134 F.3d 1400 (9th Cir. 1998), agreeing that PRWORA did not cure the constitutional defect and that respondents demonstrated a likelihood of success on the merits and irreparable harm (\href{https://www.law.cornell.edu/supremecourt/text/98-97}{Cornell LII}).
\item \textbf{U.S. Supreme Court:} Certiorari was granted, 524 U.S. 982 (1998). The case was argued January 13, 1999, and decided May 17, 1999, affirming the Ninth Circuit (\href{https://www.courtlistener.com/opinion/118286/saenz-v-roe/}{CourtListener}; \href{https://www.law.cornell.edu/supremecourt/text/98-97}{Cornell LII}).
\end{itemize}

\section{Judicial Votes}

The Supreme Court affirmed the judgment of the Ninth Circuit by a \textbf{7--2} vote (\href{https://www.law.cornell.edu/supremecourt/text/98-97}{Cornell LII}):
\begin{itemize}
\item \textbf{Majority (7):} Justice John Paul Stevens authored the opinion of the Court, joined in full by Justices Sandra Day O'Connor, Antonin Scalia, Anthony M. Kennedy, David H. Souter, Ruth Bader Ginsburg, and Stephen G. Breyer.
\item \textbf{Concurrences:} None.
\item \textbf{Dissents (2):} Chief Justice William H. Rehnquist authored a dissenting opinion joined by Justice Clarence Thomas (\href{https://www.law.cornell.edu/supct/html/98-97.ZD.html}{Cornell LII}). Justice Clarence Thomas authored a separate dissenting opinion joined by Chief Justice William H. Rehnquist (\href{https://www.law.cornell.edu/supct/html/98-97.ZD1.html}{Cornell LII}).
\end{itemize}

\section{Holding}

\textbf{Judgment:} Affirmed. California's one-year durational residency cap on welfare benefits (Cal. Welf. \& Inst. Code Ann. \S~11450.03) is unconstitutional, and Congress's authorization in 42 U.S.C. \S~604(c) cannot validate it (\href{https://www.law.cornell.edu/supremecourt/text/98-97}{Cornell LII}).

\textbf{Rule Announced:} The constitutional right to travel encompasses three components: (1) the right of a citizen of one state to enter and leave another state; (2) the right under Article IV, \S~2's Privileges and Immunities Clause to be treated as a ``welcome visitor rather than an unfriendly alien'' when temporarily present in another state; and (3) the right under the Fourteenth Amendment's Citizenship and Privileges or Immunities Clauses, for travelers who elect to become permanent residents, to be treated equally with other citizens of their new state (\href{https://www.law.cornell.edu/supremecourt/text/98-97}{Cornell LII}; \href{https://www.law.cornell.edu/constitution-conan/article-4/section-2/clause-1/right-to-travel-and-privileges-and-immunities-clause}{Constitution Annotated}). State classifications that discriminate against newly arrived bona fide citizens based on length of residence or prior state of residence are subject to a standard ``surely no less strict'' than strict scrutiny (and possibly more categorical), cannot be justified by a state's fiscal interest in saving money or deterring indigent migration, and cannot be authorized by Congress under Article I or Section 5 of the Fourteenth Amendment (\href{https://www.law.cornell.edu/supremecourt/text/98-97}{Cornell LII}).

\textbf{Holding vs. Dicta:} The \emph{holding} is that a state violates the Fourteenth Amendment's Privileges or Immunities and Citizenship Clauses when it limits newly arrived bona fide residents' welfare benefits to the level of their prior states during their first year of residence, and that federal statutory approval under PRWORA does not insulate such state discrimination from invalidation (\href{https://www.law.cornell.edu/supremecourt/text/98-97}{Cornell LII}). By contrast, the Court's discussion of the textual source of the first component of the right to travel, its observations distinguishing ``readily portable'' benefits such as divorce decrees (\textit{Sosna v. Iowa}, 419 U.S. 393 (1975)) and in-state college tuition (\textit{Vlandis v. Kline}, 412 U.S. 441 (1973)), and its reservation of what weight length of residence might carry if a newcomer's bona fide claim to state citizenship were actually contested are \emph{dicta} (\href{https://www.law.cornell.edu/supremecourt/text/98-97}{Cornell LII}).

\section{Analysis of Opinions}

\subsection{Majority Opinion (Justice Stevens)}

Justice Stevens organized the constitutional right to travel into three distinct components (\href{https://www.law.cornell.edu/supremecourt/text/98-97}{Cornell LII}). First, the right to enter and leave another state---vindicated in \textit{Edwards v. California}, 314 U.S. 160 (1941), and \textit{United States v. Guest}, 383 U.S. 745 (1966)---was not directly impaired because \S~11450.03 erected no physical barrier to entering California (\href{https://www.law.cornell.edu/supremecourt/text/98-97}{Cornell LII}). Second, Article IV, \S~2's Privileges and Immunities Clause protects temporary visitors against ``the disabilities of alienage'' unless the state has a substantial reason for distinguishing nonresidents beyond their noncitizenship, as in recreational hunting licenses or out-of-state tuition; however, those resident-versus-nonresident justifications do not apply once a citizen settles in a new state (\href{https://www.law.cornell.edu/supremecourt/text/98-97}{Cornell LII}; \href{https://www.law.cornell.edu/constitution-conan/article-4/section-2/clause-1/right-to-travel-and-privileges-and-immunities-clause}{Constitution Annotated}).

Turning to the third component, Stevens grounded the right of newly arrived permanent residents to equal treatment with established citizens squarely in Section 1 of the Fourteenth Amendment---specifically the Citizenship Clause and the Privileges or Immunities Clause (\href{https://www.law.cornell.edu/supremecourt/text/98-97}{Cornell LII}). Remarkably, Stevens noted that both Justice Miller's majority and Justice Bradley's dissent in the \href{https://www.law.cornell.edu/supremecourt/text/83/36}{\textit{Slaughter-House Cases}}, 83 U.S. (16 Wall.) 36, 80, 112--113 (1873), agreed that a U.S. citizen may choose to become a citizen of any state by bona fide residence with the same rights as other citizens of that state (\href{https://www.law.cornell.edu/supremecourt/text/98-97}{Cornell LII}). Because the discriminatory classification itself constitutes a penalty on this right, neither rational basis nor the Solicitor General's proposed intermediate scrutiny applied; the standard is ``no less strict'' than \href{https://supreme.justia.com/cases/federal/us/394/618/}{\textit{Shapiro}}'s compelling-interest test (\href{https://www.law.cornell.edu/supremecourt/text/98-97}{Cornell LII}).

Applying that standard, Stevens dismantled California's defenses. Deterring poor migrants is constitutionally impermissible and empirically unsupported (\href{https://www.law.cornell.edu/supremecourt/text/98-97}{Cornell LII}). California's fiscal rationale---saving \$10.9 million annually, which an across-the-board reduction of 72 cents per month would equally achieve---could not justify creating a hierarchy of 45 subclasses of citizens defined by length of residence and former state, because the Citizenship Clause ``does not allow for, degrees of citizenship based on length of residence'' (\href{https://supreme.justia.com/cases/federal/us/457/55/}{\textit{Zobel}}, 457 U.S. at 69) (\href{https://www.law.cornell.edu/supremecourt/text/98-97}{Cornell LII}). Unlike ``readily portable'' benefits such as college tuition or a divorce, cash welfare benefits are consumed within California, eliminating any realistic risk of sham residency (\href{https://www.law.cornell.edu/supremecourt/text/98-97}{Cornell LII}). Finally, Congress's authorization in PRWORA could not save \S~11450.03: neither Article I nor Section 5 of the Fourteenth Amendment empowers Congress to authorize states to violate Fourteenth Amendment guarantees, and the Citizenship Clause binds the federal government as well as the states (\href{https://www.law.cornell.edu/supremecourt/text/98-97}{Cornell LII}).

\subsection{Dissenting Opinion (Chief Justice Rehnquist)}

Chief Justice Rehnquist, joined by Justice Thomas, criticized the majority for reviving the long-dormant Privileges or Immunities Clause---previously used to strike down a law only once, in \textit{Colgate v. Harvey}, 296 U.S. 404 (1935), overruled five years later by \textit{Madden v. Kentucky}, 309 U.S. 83 (1940) (\href{https://www.law.cornell.edu/supct/html/98-97.ZD.html}{Cornell LII}). Rehnquist argued that the majority improperly conflated the right to travel with the distinct right to become a state citizen: once a person settles permanently in a new state, ``he is no longer `traveling' in any sense of the word'' (\href{https://www.law.cornell.edu/supct/html/98-97.ZD.html}{Cornell LII}).

Substantively, Rehnquist maintained that \S~11450.03 was a valid, objective ``good-faith residency requirement'' designed to test a newcomer's subjective intent to remain, analogous to the one-year durational residency requirements upheld for in-state university tuition (\textit{Starns v. Malkerson}, 401 U.S. 985 (1971); \textit{Vlandis v. Kline}, 412 U.S. 441 (1973)) and divorce jurisdiction (\textit{Sosna v. Iowa}, 419 U.S. 393 (1975)) (\href{https://www.law.cornell.edu/supct/html/98-97.ZD.html}{Cornell LII}). He mocked the majority's ``you can't take it with you'' portability distinction as ``border[ing] on the metaphysical,'' arguing that welfare payments invest in human capital (job training and subsistence) just as tuition subsidies do (\href{https://www.law.cornell.edu/supct/html/98-97.ZD.html}{Cornell LII}). Moreover, because California still paid newcomers their prior state's grant alongside Medicaid, job training, and increased food stamps, and because Congress expressly endorsed state flexibility in 42 U.S.C. \S~604(c), Rehnquist viewed the statute as a reasonable accommodation of state fiscal interests (\href{https://www.law.cornell.edu/supct/html/98-97.ZD.html}{Cornell LII}).

\subsection{Dissenting Opinion (Justice Thomas)}

Justice Thomas, joined by Chief Justice Rehnquist, joined Rehnquist's dissent in full and wrote separately to challenge the majority's ahistorical use of the Privileges or Immunities Clause (\href{https://www.law.cornell.edu/supct/html/98-97.ZD1.html}{Cornell LII}). Tracing the phrase ``privileges and immunities'' from the 1606 Charter of Virginia through the Articles of Confederation, Article IV, \S~2, and Justice Bushrod Washington's circuit opinion in \textit{Corfield v. Coryell}, 6 F. Cas. 546 (C.C.E.D. Pa. 1825), Thomas argued that the Thirty-Ninth Congress understood ``privileges or immunities'' to encompass only \emph{fundamental rights} belonging to citizens of all free governments---not every positive-law public benefit a state chooses to distribute, such as welfare payments (\href{https://www.law.cornell.edu/supct/html/98-97.ZD1.html}{Cornell LII}). While signaling openness to reconsidering the \href{https://www.law.cornell.edu/supremecourt/text/83/36}{\textit{Slaughter-House Cases}} in an appropriate case, Thomas warned that reviving the Clause without first ascertaining its original meaning---and without deciding whether it should ``displace, rather than augment,'' substantive due process and equal protection---risked turning it into ``yet another convenient tool for inventing new rights'' (\href{https://www.law.cornell.edu/supct/html/98-97.ZD1.html}{Cornell LII}).

\section{Examples: Future Applications}

\begin{enumerate}
\item \textbf{Same-Side (Unconstitutional):} A state enacts a law providing a \$500 monthly housing-assistance stipend to low-income residents, but limits residents who moved into the state within the past eighteen months to the average housing stipend paid in their prior state. \emph{Why:} Like California's \S~11450.03, this creates subclasses of bona fide state citizens based on duration of residence and prior domicile for a non-portable subsistence benefit consumed in-state, violating the third component of the right to travel under the Fourteenth Amendment's Privileges or Immunities Clause (\href{https://www.law.cornell.edu/supremecourt/text/98-97}{Cornell LII}).
\item \textbf{Same-Side (Unconstitutional):} A state creates a permanent homestead property-tax exemption that increases by \$1,000 for each decade a homeowner has been a continuous resident of the state, expressly rewarding ``long-term civic loyalty.'' \emph{Why:} Under \textit{Saenz} and \href{https://supreme.justia.com/cases/federal/us/457/55/}{\textit{Zobel v. Williams}}, the Citizenship and Privileges or Immunities Clauses forbid creating fixed hierarchies or degrees of state citizenship based on length of residence or past contributions (\href{https://www.law.cornell.edu/supremecourt/text/98-97}{Cornell LII}; \href{https://www.law.cornell.edu/constitution-conan/amendment-14/section-1/interstate-travel}{Constitution Annotated}).
\item \textbf{Opposite-Side (Constitutional):} A state charges out-of-state visitors \$400 for a recreational elk-hunting license while charging state residents \$40, where the visitors plan to return home after hunting season. \emph{Why:} Temporary visitors fall under the second component (Article IV, \S~2), not the Fourteenth Amendment's third component; as \textit{Saenz} reaffirmed (citing \textit{Baldwin v. Fish \& Game Comm'n of Mont.}, 436 U.S. 371 (1978)), recreational hunting is not a fundamental privilege under Article IV, and differential fees for nonresidents are permissible (\href{https://www.law.cornell.edu/supremecourt/text/98-97}{Cornell LII}).
\item \textbf{Opposite-Side (Constitutional):} A state requires that an individual reside in the state for one year before filing for divorce in state court. \emph{Why:} \textit{Saenz} expressly distinguished \textit{Sosna v. Iowa}, 419 U.S. 393 (1975), noting that a divorce decree is a ``readily portable'' legal status that could encourage sham temporary residency to obtain a decree and immediately return to one's prior state, giving the state a compelling interest in verifying bona fide domicile (\href{https://www.law.cornell.edu/supremecourt/text/98-97}{Cornell LII}).
\item \textbf{Fence-Sitter (Close Case):} A state university system requires all newly arrived students---including financially independent adults who move to the state, register to vote, pay state income taxes, and work full-time---to live in the state for two full years before qualifying for in-state tuition. \emph{Why:} On one hand, \textit{Saenz} cited \textit{Vlandis v. Kline} in dicta to distinguish college education as a ``readily portable'' benefit, and lower courts have historically upheld one-year tuition residency rules (\href{https://www.law.cornell.edu/supremecourt/text/98-97}{Cornell LII}; \href{https://repository.uclawsf.edu/cgi/viewcontent.cgi?article=3500&context=hastings_law_journal}{Hastings L.J.}). On the other hand, a rigid two-year waiting period applied to demonstrably bona fide permanent residents who have already established domicile goes beyond testing intent to remain and risks functioning as an impermissible penalty on new state citizens under \textit{Saenz}'s strict categorical rule.
\end{enumerate}

\section{Critique}

\textbf{Originalist and Textualist Perspectives:} Originalist scholars have greeted \textit{Saenz} with sharply divided reactions. Some praised the Court for breaking its 126-year silence on the Fourteenth Amendment's Privileges or Immunities Clause after the \href{https://www.law.cornell.edu/supremecourt/text/83/36}{\textit{Slaughter-House Cases}} gutted the text (\href{https://www.cato.org/sites/cato.org/files/2020-04/shapiro-blackman-gmlr-priv-imm-clause.pdf}{Shapiro \& Blackman, Cato/GMU L. Rev.}; \href{https://lawliberty.org/forum/reversing-the-legacy-of-slaughter-house/}{Wurman, Law \& Liberty}). Yet many originalists echo Justice Thomas's critique: Stevens invoked the Clause without engaging its Reconstruction-era history or reconciling it with \textit{Slaughter-House}, and stretching ``privileges or immunities'' to positive-law welfare entitlements sits uneasily with the nineteenth-century \textit{Corfield} tradition of fundamental civil rights (\href{https://www.law.cornell.edu/supct/html/98-97.ZD1.html}{Cornell LII}; \href{https://journals.law.harvard.edu/jlpp/wp-content/uploads/sites/90/2026/03/30_23HarvJLPubPoly2951999-2000.pdf}{Harvard JLPP}). Moreover, scholars like Kurt Lash and John Harrison debate whether the Clause primarily protects enumerated national rights or serves as an equality/antidiscrimination guarantee among state citizens (\href{https://lawliberty.org/forum/reversing-the-legacy-of-slaughter-house/}{Wurman, Law \& Liberty}; \href{https://journals.law.harvard.edu/jlpp/wp-content/uploads/sites/90/2026/03/30_23HarvJLPubPoly2951999-2000.pdf}{Harvard JLPP}).

\textbf{Progressive and Doctrinal Perspectives:} Progressive commentators welcomed \textit{Saenz}'s categorical protection of poor migrants against two-tier welfare schemes and its affirmation that Congress cannot authorize Fourteenth Amendment violations (\href{https://www.legalmomentum.org/library/saenz-v-roe}{Legal Momentum}). At the same time, doctrinal scholars---notably Laurence H. Tribe in \textit{Saenz Sans Prophecy} and Linda Bosniak---warned that shifting rights from the Equal Protection Clause (which protects ``any person'') to the Privileges or Immunities Clause (which protects ``citizens of the United States'') could inadvertently weaken constitutional protections for noncitizens and lawful permanent residents (\href{https://kb.osu.edu/bitstream/handle/1811/70511/OSLJ_V63N5_1285.pdf}{Ohio St. L.J.}). Additionally, critics noted that \textit{Saenz}'s revival of the Clause proved largely a one-off: in \textit{McDonald v. City of Chicago}, 561 U.S. 742 (2010), only Justice Thomas relied on Privileges or Immunities, while the plurality returned to substantive due process (\href{https://www.cato.org/sites/cato.org/files/2020-04/shapiro-blackman-gmlr-priv-imm-clause.pdf}{Shapiro \& Blackman, Cato/GMU L. Rev.}).

\textbf{Underlying Competing Values:} The decision reflects a deep structural tension between \emph{national union and equal citizenship}---the premise that Americans form one Nation in which citizens choose their states rather than states selecting their citizens---and \emph{fiscal federalism}, which seeks to allow states to experiment with generous social safety nets without fearing a ``race to the bottom'' triggered by interstate benefit disparities (\href{https://www.law.cornell.edu/supremecourt/text/98-97}{Cornell LII}; \href{https://law.lclark.edu/live/files/9702-schwinnpdf}{Schwinn, Lewis \& Clark L. Rev.}).

\section{Key Quotations}

\begin{quote}
``The `right to travel' discussed in our cases embraces at least three different components. It protects the right of a citizen of one State to enter and to leave another State, the right to be treated as a welcome visitor rather than an unfriendly alien when temporarily present in the second State, and, for those travelers who elect to become permanent residents, the right to be treated like other citizens of that State.'' \\
--- Justice Stevens, majority opinion (\href{https://www.law.cornell.edu/supremecourt/text/98-97}{526 U.S. at 500, Cornell LII})
\end{quote}

\begin{quote}
``Citizens of the United States, whether rich or poor, have the right to choose to be citizens `of the State wherein they reside.' U. S. Const., Amdt. 14, \S~1. The States, however, do not have any right to select their citizens.'' \\
--- Justice Stevens, majority opinion (\href{https://www.law.cornell.edu/supremecourt/text/98-97}{526 U.S. at 510--511, Cornell LII})
\end{quote}

\begin{quote}
``Congress has no affirmative power to authorize the States to violate the Fourteenth Amendment and is implicitly prohibited from passing legislation that purports to validate any such violation.'' \\
--- Justice Stevens, majority opinion (\href{https://www.law.cornell.edu/supremecourt/text/98-97}{526 U.S. at 508, Cornell LII})
\end{quote}

\begin{quote}
``A person is no longer `traveling' in any sense of the word when he finishes his journey to a State which he plans to make his home.'' \\
--- Chief Justice Rehnquist, dissenting (\href{https://www.law.cornell.edu/supct/html/98-97.ZD.html}{526 U.S. at 513, Cornell LII})
\end{quote}

\begin{quote}
``We should also consider whether the Clause should displace, rather than augment, portions of our equal protection and substantive due process jurisprudence.'' \\
--- Justice Thomas, dissenting (\href{https://www.law.cornell.edu/supct/html/98-97.ZD1.html}{526 U.S. at 528, Cornell LII})
\end{quote}

\end{document}
